\documentclass[../main.tex]{subfiles}
\begin{document}
%==========================================TIÊU ĐỀ TẬP========================================
\begin{table}[h]
  \begin{adjustwidth}{-1cm}{}
    \begin{tabular}{>{\centering\arraybackslash}p{6cm}|>{\raggedright\arraybackslash}p{12.5cm}}
      \multirow{5}{6cm}
      % Cột bên trái: ÉPISODE và số tập
      {\\[-40pt]
        \flaregothic\fontsize{20pt}{20pt}\selectfont \centering
        \color{eptype}ÉPISODE\color{black}\\
        \burbank\fontsize{72pt}{72pt}\selectfont
        \color{epnum}58\color{black}\\[6pt]
        \cafeteria\fontsize{20pt}{20pt}\selectfont\color{epnum}(5-2)\color{black}
        \\[18pt]
        \flaregothic\fontsize{14pt}{14pt}\selectfont
        \color{epattr}BIENVENUE, \uline{\color{Shiranui} Flare} !
        \color{black}
        \\}
       &
      % Dòng thứ nhất: tên tập tiếng Nhật
      {\fotskip\fontsize{18pt}{18pt}\selectfont くさや\ruby{食}{た}べてみた【\ruby{臭}{くさ}い】
      } \\[6pt]
      % Dòng thứ hai: tên tập tiếng Anh
       & {\cafeteria\fontsize{18pt}{18pt}\selectfont Kusaya Eater} \\[4pt]
      % Dòng thứ ba: tên tập tiếng Việt
       & {\vnmsans\fontsize{18pt}{18pt}\selectfont Món cá khô bốc mùi} \\[8pt]
      % Dòng thứ tư: Nhân vật xuất hiện
       & {\caviar{\fontsize{14pt}{14pt}\selectfont Nhân vật: \emp{kuon1} \emp{pekora} \emp{rushia}\emp{flare1}  \emp{noel} \emp{marine} }}
      \\[8pt]
      % Dòng cuối cùng: Ngày phát hành
       & {\continuum\fontsize{16pt}{16pt}\selectfont Ngày phát hành: 14/06/2020} \\[18pt]
    \end{tabular}
  \end{adjustwidth}
\end{table}
%=========================================TRUYỆN CHÍNH========================================
\color{black}

% Hộp tóm tắt
\txtbx[2]{{\color{vol} \uline{Tóm tắt:}} Marine mang một miếng cá đuối khô (kusuya) vào văn phòng, gây ra thảm họa mùi hôi cực độ. Flare - tân binh của nhóm Fantasy - cùng Kuon, Rushia và Noel tìm cách khử mùi nhưng bất thành. Pekora bị Rushia bỏ vào hộp ``xin nhận nuôi'' và mất củ cà rốt. Sau khi cuốn cá khô trôi đi theo nước mắt lụt lội của Pekora, cả nhóm phát hiện thủ phạm thực sự của mùi hôi thứ hai là\dots\ một lọ mắm tôm để quên hàng tháng của Kuon. Kết quả: thế hệ thứ Ba gục ngã hoàn toàn.}

Mùa hè đã thực sự gõ cửa. Cái nóng tuy chưa đến mức gay gắt, nhưng cũng đủ để những tia nắng vàng óng trải dài trên khung cửa sổ, báo hiệu một mùa oi ả sắp bủa vây. Thế nhưng, điều đặc biệt hôm nay không phải là nhiệt độ, mà là\dots\ một mùi hương vô cùng đặc biệt. Một thứ mùi mà chỉ cần hít một hơi là biết ngay: không phải nước hoa, cũng chẳng phải đồ ăn thơm phức. Là mùi hôi. Hôi một cách kỳ lạ, nồng nặc, và dường như đang lan tỏa khắp mọi ngóc ngách của văn phòng \hll.

\chrint

\model{23}{r}{6cm}{flare1}

Trong bầu không khí ngột ngạt ấy, một thành viên mới xuất hiện. Shiranui Flare (\ruby{不知火}{しらぬい}フレア) - nàng Yêu tinh vàng, thành viên cuối cùng của Thế hệ thứ Ba. Cô là một con lai giữa người và yêu tinh, mang trong mình một tâm hồn thích tận hưởng cuộc sống theo cách riêng, nhưng cũng rất dễ\dots\ mệt mỏi. Flare sống theo cảm xúc, đa cảm, dễ khóc và dễ xấu hổ, nhưng nếu được khen, cô sẽ vui cả ngày, dù có hơi ngượng. Bên cạnh cô luôn có một tiểu tiên tên Kintsuba, một sinh vật nhỏ nhắn trông chẳng khác gì gấu trúc con.

Flare tốt bụng, cực kỳ khiêm tốn, đôi khi khiêm tốn đến mức thành khuyết điểm. Lòng tốt của cô được cả nhóm Fantasy chứng thực, và điều đó thường dẫn đến một ``harem'' không chủ đích xung quanh cô, ngoại trừ Pekora. Trong số bốn thành viên thế hệ Ba, Noel là người thể hiện tình cảm rõ ràng nhất.

Các thành viên khác như Sora cũng từng bày tỏ sự ngưỡng mộ dành cho Flare, vì giọng nói đa dạng và khí chất trưởng thành. Dù nhìn qua có vẻ lạnh lùng, xa cách, Flare thực chất khá trẻ con và nữ tính, điều này lộ rõ nhất khi cô chơi game kinh dị hay khi say rượu.

Flare là một trong số ít thành viên mà Kuon gọi là ``học sinh gương mẫu'' (chỉ sau Sora và Mio). Cô nghiêm túc, đôi khi ngán ngẩm trước những trò điên rồ của đồng nghiệp, và sẵn sàng hỗ trợ cậu Tổng dọn dẹp mớ rắc rối do các cô nàng \hll\ bày ra.

Về ngoại hình, Flare sở hữu làn da ngăm đen cùng đôi tai nhọn đặc trưng của yêu tinh. Mái tóc vàng hoe dài đến giữa lưng, thường được buộc đuôi ngựa cao bằng nơ xanh đậm, tóc tết hai bên. Trước trán có những sợi trắng. Cô có một bộ ngực khá lớn.

Trang phục của cô là sự chắp vá từ nhiều lớp vải: bên trong là áo trắng không dây, lớp áo nịt ngực bên ngoài, cổ áo và áo choàng xanh nước biển đậm có đường cắt ở vai và lưng trên, hai tay áo rời màu xanh đậm cắt ở khuỷu tay. Phía dưới là quần soóc trắng, tất trắng với viền đùi ca-rô xanh-trắng ở chân phải, giày cao gót xanh-trắng. Cô đeo găng tay đen xỏ ngón trỏ.

\chrint

Flare vừa bước vào, đôi tai nhọn khẽ động đậy như một radar đang quét tìm nguồn gây nhiễu. Cô nhíu mày, đưa tay quạt quạt trước mũi: ``Sao hôm nay trong này hôi thế nhỉ? Mùi này\dots\ không phải đồ ăn thiu, cũng chẳng phải ẩm mốc. Nó kinh khủng theo một cách rất riêng.''

Kuon ngồi ở bàn làm việc, ngước lên nhìn quanh. Cậu thở hắt ra, giọng bình thản: ``Em thấy thế nào chứ anh thấy bình thường mà. Có khi do em mới đến nên chưa quen không khí ở đây. Văn phòng lúc nào chẳng có mùi\dots\ lai rai.''

Noel đang ngồi ở ngay bên cạnh nàng Yêu tinh, tay bịt mũi nãy giờ. Cô lên tiếng, giọng nghèn nghẹt như đang bị cảm lạnh: ``Nãy giờ đang cố gắng nín thở nên mình chẳng nhận ra luôn. Có gì đâu.''

\img{5-2a}

Flare liếc mắt sang cô bạn, đôi lông mày nhướng lên: ``Nín làm chi vậy?''

Noel bỏ một tay ra khỏi mũi, nhưng tay kia vẫn bịt chặt. Cô hít một hơi thật sâu bằng miệng (trông như một vận động viên đang tập thở dưới nước), rồi đáp: ``Chả là sắp tới hè rồi nên mình đang cố gắng để có thể ở dưới nước lâu nhất có thể\dots\ Tập thở, tập nín, luyện phổi.''

Cả Flare và Kuon đồng thanh nhíu mày, giọng đầy ngờ vực: ``Thì sao?''

Nữ Hiệp sĩ đập tay xuống bàn một cái rầm, mắt sáng rực, đầy quyết tâm: ``Để trở thành một VTuber hàng đầu còn gì nữa! Ai bảo VTuber chỉ ngồi trước màn hình? Phải đa năng, phải biết bơi lội, phải chinh phục cả đại dương chứ!''

Cả căn phòng chìm vào một khoảng lặng dài, ngập tràn sự bối rối. Cậu Tổng nhìn Noel như thể cô vừa tuyên bố rằng mình sẽ lên mặt trăng bằng xe đạp.

Nàng Yêu tinh vàng khoanh tay, thở dài: ``Tôi nghĩ là cậu nên tìm cách khác thì hay hơn. Luyện nín thở để ở dưới nước lâu\dots\ rồi ra sao? Livestream dưới hồ bơi à?''

Cậu Tổng lắc đầu ngao ngán: ``Anh thấy cái trò này ngớ ngẩn thế nào ấy. Kẻo rước họa vào thân đấy. Có khi đang livestream thì ngất xỉu vì thiếu oxy, khán giả tưởng em lại đang biểu diễn nghệ thuật thì khổ.''

Flare gật gù đồng tình, nhưng rồi cô lại cau mày, đôi mắt liếc quanh phòng. Vấn đề lớn hơn vẫn chưa được giải quyết: mùi hôi vẫn còn đó, nồng nặc, dai dẳng, như một linh hồn ma quái không chịu rời đi.

``Mà\dots\ cái mùi hôi thối này bắt nguồn từ đâu nhỉ? Không phải từ Noel, cũng chẳng phải từ đồ dùng văn phòng. Hình như\dots\ từ trong bếp.''

Đúng lúc này, trong căn bếp nhỏ của văn phòng, Marine và Rushia đang đứng như hai pho tượng trước một thứ vô cùng đặc biệt. Trên thớt, một miếng cá đuối khô (kusuya) nằm chễm chệ, da nâu sậm, mình cứng queo, và tỏa ra một thứ mùi mà có lẽ cả người chết cũng phải sống dậy để\dots\ bịt mũi.

\img{5-2b}

Nàng Thuyền trưởng chỉ tay vào miếng cá, giọng đầy tự hào: ``Đây là con cá khô kusuya mà tôi chuẩn bị cho buổi stream sắp tới đây. Hàng đặc biệt, nhập từ\dots\ ờm, từ một nơi nào đó có mùi mạnh.''

Nàng Pháp sư đứng bên cạnh, tay bịt mũi, mắt nhắm nghiền, giọng dứt khoát như một bản án: ``Vứt nó đi.''

Marine giật mình, quay sang: ``Này, đừng có lãng phí đồ ăn như thế chứ! Cá khô là cả một nghệ thuật, bà biết không?''

Rushia không trả lời. Cô chỉ càng lúc càng lùi xa, mắt nhìn miếng cá với vẻ khinh miệt pha lẫn sợ hãi.

Marine không bỏ cuộc. Cô sán lại gần, giọng bỗng trở nên\dots\ ngọt ngào một cách đáng sợ: ``Vậy là hư lắm nhá\~{} Của hiếm đấy, không dễ gì mua được.'' Cô vừa nói vừa đưa tay chỉ chỉ, mắt híp lại, tạo nên một biểu cảm vừa dâm đãng vừa đe dọa, kiểu ``mày không ăn thì tao cho mày ăn đấm''.

Rushia tỏ vẻ khó chịu tột độ. Cô hét lên, giọng cao hơn bình thường: ``Vậy thì bà đem về mà ăn đi, chứ mang tới đây làm gì!? Mở nắp ra là cả tòa nhà sập rồi!''

Thấy đồng nghiệp mắng mình to tiếng, Senchou giật mình. Mắt cô mở to, lắp bắp: ``R-Rushia? Có cần phải thế không? Chỉ là một miếng cá thôi mà\dots''

Nhưng Rushia chẳng thèm nán lại. Cô quay lưng chạy thẳng ra cửa, vừa chạy vừa hét, giọng vang dội khắp hành lang: ``\textbf{Hôi quá!!}''

\gif{5-2c}{1}{12}

Tiếng hét của nàng Pháp sư như một hồi còi báo động, lập tức thu hút sự chú ý của Kuon và Flare, khi hai người đang dò dẫm tìm nguồn cơn của mùi hôi. Họ bước vào bếp, và cảnh tượng trước mắt khiến mọi thắc mắc tan biến.

Nàng Yêu tinh vàng bịt mũi, nhíu mày, giọng nghẹn lại: ``Vậy ra mùi hôi bắt nguồn từ đây à. Tôi cứ tưởng có ai để quên xác chuột trong tủ.''

Cậu Tổng gật đầu, mặt vẫn bình thản kỳ lạ: ``Chả trách sao có mùi thum thủm thế. Mà, có chuyện gì mà Ru-chan hét lên vậy?''

Marine quay sang nhìn hai người mới đến, vẻ mặt vẫn còn bối rối sau phản ứng dữ dội của Rushia: ``C-Chắc em nghĩ là cậu ấy sợ mùi thôi. Hồi nãy còn định mời cậu ấy ăn chung, ai ngờ\dots\ Mà hai người muốn ăn không?''

Flare lắc đầu nguầy nguậy, giọng dứt khoát như một lời thề: ``Thôi, miễn! Tôi còn muốn sống đến cuối ngày.''

Cậu Tổng nhẹ nhàng hơn, nhưng cũng chẳng kém phần kiên quyết: ``Ờm, thôi. Anh không thích ăn cá lắm đâu.''

Senchou hơi nhướng mày, tò mò: ``Anh không thấy nó bốc mùi ghê sao? Hay anh bị nghẹt mũi?''

Cậu chỉ nhún vai, điềm nhiên đáp: ``Anh từng ngửi mấy thứ còn kinh hơn thế nhiều. Hồi xưa đi thực tập ở nhà máy xử lý rác thải, có những hôm gió nồm, mùi cá thối ủ chua, rồi còn mùi rác bốc lên nồng nặc nữa\dots\ So với đây thì miếng cá khô này còn dễ thở chán.''

Flare và Marine liếc nhìn nhau, một thoáng im lặng. Rồi cả hai cùng ``ồ'' một tiếng, không hỏi thêm. Có lẽ, có những thứ trong quá khứ của cậu quản lý mà tốt nhất không nên đào sâu.

\ct

Bất chợt, cánh cửa bật mở. Rushia quay trở lại, nhưng lần này cô không đến một mình. Trên tay, nàng Pháp sư ôm chặt Pekora đang run lẩy bẩy, mắt nhắm nghiền, trông chẳng khác gì một chú thỏ con bị ai đó vớt ra khỏi chuồng giữa đêm mưa bão.

Trên lọn tóc của Pekora, có một nhánh gừng và một cây hành lá cài lên một cách lộn xộn, như thể ai đó đã cố gắng trang trí cô như một món ăn. Thực ra, chúng là những thứ tạm bợ thay cho vài củ cà rốt đã bị ai đó cuỗm mất để làm món cà ri\footnote{Xem lại {\flaregothic ÉPISODE 45}, phần {\abv Truyện phụ}, quyển số Bốn.}.

Nàng Yêu tinh vàng nhướng mày, giọng đầy ngạc nhiên: ``Đù, bà ấy quay lại nhanh thế. Mới chạy ra được vài phút đã kéo thêm người về. Tốc độ ánh sáng à?''

Cậu Tổng cũng tròn mắt, nhưng chưa kịp hỏi thì Rushia đã thở hổn hển, vừa nói vừa rung rung Pekora như một con búp bê vải: ``Tôi nghe nói là gừng với hành tây có thể khử mùi hôi đấy! Mà tôi không có hành tây, nên đành dùng hành lá tạm. Với lại, tôi cần một cái\dots\ vật chứa để gắn chúng vào.''

Cậu Tổng gật gù, nhưng vẻ mặt vẫn đầy hoài nghi: ``Hoặc là tinh dầu, hoặc là giấm, hoặc là than hoạt tính\dots\ Nói chung còn nhiều cách khử mùi mà chẳng cần đến gừng hay hành. Cần quái gì phải mang em ấy sang đây? Em ấy có phải là máy lọc không khí đâu.''

Nghe đến đó, nàng Thỏ mở mắt lờ đờ, giọng ngái ngủ, như thể vừa tỉnh dậy sau một cơn mê dài: ``Thế\dots\ bà vác tôi tới đây làm gì hả, peko\dots? Tôi đang mơ thấy mình đang ăn cà rốt ở thiên đường thì bà bỗng nhiên xách dậy.''

\img{5-2d}

Rushia nhìn xuống Pekora, nghĩ rằng câu hỏi cũng hợp lý. Cô đáp, giọng đầy bực bội pha chút xấu hổ: ``Vậy để tôi đưa bà về nơi tôi đã nhặt vậy! Khỏi phiền!''

Cánh cửa đóng sầm lại. Tiếng *RẦM!* vang lên khiến cả Flare và Marine đều giật mình.

Nàng Yêu tinh vàng nhìn về phía cánh cửa, rồi quay sang cậu Tổng: ``Nhặt về? Ý cậu ấy là sao chứ? Pekora-chan bị nhặt ở đâu? Trong thùng rác à?''

Cậu Tổng chỉ nhún vai, giọng bất lực: ``Anh cũng chẳng biết nữa. Chắc là tìm thấy em ấy đang loanh quanh ở đâu rồi tha về đây mà.''

Hóa ra, sau khi bước ra ngoài, Rushia không hề đặt Pekora xuống đất. Cô thản nhiên bế nàng Thỏ đi thẳng đến một góc của thành phố, nơi có một chiếc hộp giấy lớn nằm sẵn từ lúc nào không rõ. Cô nhẹ nhàng thả Pekora vào trong hộp như đang đặt một con mèo hoang.

Trên thành hộp, có một dòng chữ nguệch ngoạc, viết bằng bút dạ đỏ: ``\uline{Xin hãy nhận nuôi bé.}''

\img{5-2e}

Pekora ngồi trong hộp, mắt mở to, chớp chớp vài lần. Cô nhìn xung quanh với bốn bề là con đường vắng tanh, bầu trời và những tòa nhà\dots\ rồi nhìn xuống dòng chữ ``xin hãy nhận nuôi''. Rồi cô nhìn lên Rushia, người đang đứng ngoài hộp với vẻ mặt vô cùng hài lòng.

``\textbf{Khoan, có gì đó sai sai ở đây, peko!!}'' Pekora hét lên, giọng đầy hoảng loạn. ``Bà định bỏ rơi tôi à!? Tôi không phải mèo hoang, peko!''

Rushia không đáp. Cô chỉ nở một nụ cười ma mãnh, đôi mắt híp lại, và\dots

Nàng Thỏ vẫn đang ngồi trong hộp, hai tay bám vào thành, cố gắng trèo ra. Nhưng bỗng nhiên, cô khựng lại. Cô cảm thấy có gì đó thiếu thiếu trên đỉnh đầu. Hơi nhẹ hơn, hơi trống trải hơn.

Cô đưa tay lên sờ, nơi vốn có một củ cà rốt nhỏ xinh cắm trên tóc, và cảm nhận được\dots\ không khí.

Củ cà rốt cuối cùng cô để dành cho buổi livestream chiều nay đã không cánh mà bay!

Pekora ngước mắt lên, và kịp nhìn thấy cô bạn đang cầm củ cà rốt trên tay, hai ngón tay xoay xoay như một trái bóng nhỏ. Trên môi nàng Pháp sư là nụ cười đắc thắng.

``Hình như cái này cũng có thể khử mùi đấy\~{}'' Rushia ngân nga, giọng vô cùng thích thú, ``Gừng, hành, cà rốt - combo ba món, đảm bảo sạch mùi luôn nanodesu\~{}''

\img{5-2f}

Pekora thấy thế, mắt đỏ hoe, cô vùng vẫy trong hộp, hét lên đầy tuyệt vọng: ``\textbf{Trả nó lại cho tôi ngay, peko!!} Đó là củ cà rốt cuối cùng tôi để dành! Tôi đã chăm bón cho nó suốt ba tháng, peko!''

Nhưng Rushia chẳng thèm để ý. Cô vừa tung hứng củ cà rốt vừa bước ra phía cửa sổ, giọng vẫn ngân nga.

Và nàng Thỏ, ngồi trong thùng giấy với dòng chữ ``xin hãy nhận nuôi'', chỉ còn biết gào lên trong vô vọng. ``\textit{Tại sao\dots\ tại sao tôi lại bị bỏ rơi vì một miếng cá khô chứ, peko\dots}''

\ct

Quay trở lại căn bếp, không khí vẫn ngột ngạt và nồng nặc. Flare, Marine và Kuon đang đau đầu tìm cách xử lý miếng cá khô, thứ mà giờ đây trông chẳng khác gì một vũ khí hủy diệt hàng loạt. Mỗi người một vẻ mặt: nàng Thuyền trưởng thì bịt mũi nhăn nhó, mắt đỏ hoe như sắp khóc; cậu con trai chống tay bĩu môi, trông như một giáo sư đang suy nghĩ về một bài toán siêu việt; còn nàng Yêu tinh vàng lại\dots\ khá vui vẻ. Cô không quá lo lắng với mùi này, thậm chí còn hơi tò mò: ``Mùi này có thể đánh bật được cả quái vật trong game kinh dị của tôi không nhỉ?''

Flare lên tiếng đề xuất, giọng đầy phấn khích của một người thích thử nghiệm: ``Mình cứ thử bỏ hết hành tây với gừng vào rồi nấu lên xem. Theo khoa học, những thứ này có tinh dầu khử mùi mạnh lắm.''

Marine vốn đã nhận được một ít gừng và hành lá từ Rushia (sau vụ ``bắt cóc'' Pekora) bèn hí hửng bỏ tất cả vào nồi. Cô bật bếp, đợi nước sôi, và\dots

\img{5-2g}

Một làn khói đen đặc, đặc quánh, bốc lên từ chiếc nồi. Không phải khói trắng tinh khiết của những món ngon, mà là một thứ hắc khí, u ám, như thể có một con quỷ nhỏ đang bay ra từ miệng nồi. Mùi càng lúc càng nồng, càng khó chịu, khiến Marine phải lùi lại hai bước.

``Không ổn rồi, con cá thối này mạnh quá!'' cô nhăn mặt, lấy tay quạt khói. ``Nó biến hành với gừng thành tro hết rồi! Nghe này, hành và gừng đang kêu cứu kìa!''

Cậu Tổng nhìn khung cảnh trước mặt, lắc đầu: ``Xem chừng có vẻ khó rồi đấy. Con cá này không phải dạng vừa. Nó như thể đã\dots\ đắc đạo rồi ấy.''

Flare nghi ngờ nhìn con cá, đôi mắt híp lại: ``Có thật là\dots\ nó chỉ là con cá khô thối không vậy? Hay đây là một sinh vật ngoài hành tinh cải trang?''

Cậu con trai nheo mắt quan sát miếng cá. Bỗng anh giật mình: ``Trông con cá này không phải cá bình thường\dots\ Cỡ như con cá đuối chứ chả đùa. Cá đuối khô, thứ này nổi tiếng là bốc mùi nhất trong các loại cá khô. Và miếng này chắc đã được ủ mấy năm nay.''

Marine tuyệt vọng, hai tay bóp đầu: ``Ta còn cách nào để cứu vãn nữa không? Chẳng nhẽ phải vứt đi? Phí quá, em mua nó bằng cả tháng lương part-time đấy!''

Kuon xoa cằm suy nghĩ, trông như một nhà khoa học đang tìm kiếm phương án cuối cùng trước khi thảm họa xảy ra: ``Nếu mùi mạnh đến mức hành gừng còn bị đánh bại, thì anh không nghĩ giấm chua hay tinh dầu có thể hạ được nó đâu. Có khi cần đến mặt nạ phòng độc hoặc\dots\ gọi đội xử lý chất thải nguy hại.''

Trong lúc ba người đang tính kế cho miếng cá khô, nàng Thuyền trưởng để ý thấy Noel ở gần đó vẫn đang bịt mũi, mặt đỏ bừng, trông như sắp nổ tung. Marine tò mò hỏi: ``Sao đằng ấy lại bịt mũi mình vậy? Hôi quá à? Sao không chạy ra ngoài cho khỏe?''

Noel vẫn giữ vẻ mặt nghiêm túc, mắt nhìn xa xăm như một võ sĩ chuẩn bị bước vào trận đấu cuối cùng: ``Tăng thời gian tôi có thể ở dưới nước! Mỗi giây nhịn thở là mỗi giây vàng cho sự nghiệp VTuber dưới lòng đại dương của tôi!''

Cậu Tổng liếc nhìn cô nàng, giọng đầy ngán ngẩm: ``Nãy giờ em vẫn đang làm trò con bò vô nghĩa này à? Từ lúc anh vào bếp đến giờ, em đã bịt mũi được gần tiếng rồi. Gan thật.''

Flare cũng nhướng mày khó hiểu: ``Chẳng phải cậu đã nhịn thở được gần tiếng rồi à? Thành tích này đủ phá kỷ lục thế giới rồi đấy. Thôi bỏ tay ra đi, mặt cậu xanh như tàu lá rồi kìa.''

Noel lúc này bắt đầu run rẩy, đôi vai khẽ rung lên, mắt hoa lên: ``M-Mình chịu hết nổi rồi\dots\ Không thể cố thêm\dots\ chút nữa\dots\ Phổi sắp nổ tung rồi\dots''

Thấy cô nàng chuẩn bị bỏ tay ra, Kuon lập tức lên tiếng cảnh báo, giọng gấp gáp: ``Ấy, đừng có thở ở đây-''

Nhưng đã quá muộn.

Nữ hiệp sĩ, với phản xạ của một người vừa bị đuối nước trên cạn, buông tay khỏi mũi và hít một hơi thật sâu, một hơi mà có lẽ cô đã dồn nén cho cả một tuần.

Nàng Yêu tinh vàng hét lên: ``Này, đừng có hít sâu thế khi đang đứng gần con cá thối chứ!''

Và rồi\dots\ sự việc xảy ra trong tích tắc.

\img{5-2h}

Máu mũi Noel phun ra như suối; một dòng đỏ tươi, mạnh mẽ, như thể cô vừa trúng phải một đòn knock-out từ một võ sĩ hạng nặng. Cô trợn mắt lên, lảo đảo, rồi lăn đùng ra ngất xỉu, mặt vẫn còn hướng về phía chiếc nồi cá khô đầy uy lực.

Flare hoảng hốt lao tới, cố dìu lấy Noel vừa lịm đi: ``\textbf{Noel!!} Tỉnh lại đi cậu! Đừng chết ở đây, bên cạnh con cá thối! Sẽ không ai dám đến gần để đưa tang cậu đâu!''

Cậu Tổng đứng nhìn, thở dài bất lực, tay xoa trán: ``\dots\ Rồi xong. Ngu thì dính chưởng thôi. Anh đã bảo rồi, đừng có hít thở gần con cá này. Giờ thì Noel-chan thành nạn nhân đầu tiên của `thảm họa cá khô'.''

Marine đứng bên cạnh, tay vẫn bịt mũi, giọng nghẹn lại: ``Có cần gọi xe cứu thương không?''

Flare vừa lắc Noel vừa quát: ``\textbf{Gọi cứu hỏa đi! Mang bình oxy vào đây! Không, mang mặt nạ phòng độc luôn!}''

Cậu Tổng chỉ biết lắc đầu, bước lại gần hơn, kiểm tra mạch của Noel. Anh thở phào: ``Vẫn còn sống. Nhưng lần sau nhớ đừng tập thở dưới nước bằng cách nhịn thở gần cá thối nữa. Tập ở bể bơi, ở hồ, ở bồn tắm\dots\ đâu chẳng được.''

Nàng Yêu tinh ôm lấy cô bạn, vừa lo lắng vừa tức giận: ``Lát nữa tỉnh dậy, tôi sẽ bắt cậu uống một lít tinh dầu xả để khử mùi! Chết tiệt, cái công ty này đúng là hỗn loạn mà!''

\ct

Sau bao nỗ lực vô vọng, bao thí nghiệm thất bại, bao lần khói đen bốc lên như phim kinh dị, cuối cùng cả nhóm buộc phải thừa nhận một sự thật phũ phàng: miếng cá thối kia đã thực sự vô phương cứu chữa. Nó đã chiến thắng. Nó đã chứng minh được sức mạnh vượt trội so với hành, gừng, giấm, và cả\dots\ sức chịu đựng của con người.

Họ đã phải bỏ miếng cá khô ấy đi.

Flare lầm bầm trong miệng, giọng đầy bực bội nhưng cũng không kém phần ngao ngán: ``Marine thật là\dots\ Đừng có bày ra rồi cho người khác dọn chứ! Lần sau muốn ăn cá khô thì về nhà mà ăn, đừng mang vào công ty. Cả tòa nhà sắp phải sơ tán vì mùi rồi.''

\img{5-2i}

Cô vừa lẩm bẩm vừa nhấc miếng cá thối lên bằng hai đầu ngón tay, trông như thể đang cầm một quả bom hẹn giờ. Cô bước ra ngoài, định bụng đem nó vứt vào thùng rác ở cuối hành lang, hoặc một nơi nào đó càng xa càng tốt, tốt nhất là sang tận tỉnh khác.

Nhưng khi vừa ra đến hành lang, Flare liền bắt gặp một cảnh tượng khiến cô khựng lại.

Pekora vẫn đang đứng co ro trong cái hộp giấy mà Rushia đã để ngoài ban nãy. Chiếc hộp với dòng chữ ``Xin hãy nhận nuôi bé'' vẫn còn nguyên, nhưng Pekora thì chẳng còn sức để trèo ra nữa. Cô chỉ biết ngồi thu mình, hai tay ôm đầu gối, mắt thẫn thờ.

Flare tiến lại gần, đôi mắt dừng lại trên lọn tóc đặc trưng của Pekora, nơi vốn thường có một củ cà rốt nhỏ xinh cắm trên đó. Giờ thì trống trơn. Cô nhíu mày, tỏ vẻ ngạc nhiên: ``Cậu có thấy kỳ kỳ khi chẳng có gì mắc lên tóc mình không? Trông cậu như thiếu đi một phần quan trọng của bộ mặt vậy.''

Pekora ngước lên, mắt vẫn còn hoe đỏ. Cô sợ hãi thu mình lại, giọng run run: ``Ê\dots\ Đừng có nói là\dots\ bà định nhét con cá thối vào tóc tôi đấy chứ, peko? Hồi tôi lần đầu lên văn phòng, tóc tôi từng hút cả một cái tủ lạnh\footnote{Xem lại {\flaregothic ÉPISODE 34}, quyển số Ba.}, nhưng mà cá thối thì không, peko! Không bao giờ!''

Flare nhìn cô, rồi bất ngờ bật cười. Cô không nhét cá, cũng chẳng trêu đùa gì thêm. Thay vào đó, nàng Yêu tinh vàng nhẹ nhàng đặt vào lọn tóc trống trơn của Pekora\dots\ một củ cà rốt tươi rói, cam sáng, còn nguyên đất, đúng vị trí mà nó vốn nên ở.

``Đùa thôi. Cậu biết là tôi sẽ không làm thế đâu mà nhỉ?''

\img{5-2j}

Pekora nhìn củ cà rốt trên đầu mình, rồi nhìn Flare, người vẫn đang nở một nụ cười hiền. Đôi mắt cô bắt đầu rưng rưng, môi mấp máy. Cô nấc lên một tiếng, rồi đột nhiên\dots

``\textbf{Flare\~{} Waaaah!!}''

Nàng Thỏ òa khóc. Không phải khóc lóc thông thường, mà là một trận khóc nức nở, xối xả, mãnh liệt như thể cô vừa tìm lại được gia tài sau một trận hỏa hoạn. Dòng nước mắt tuôn ra như thác đổ, như vòi rồng, như một cơn mưa nhiệt đới giữa lòng thủ đô.

Và rồi, như một hiệu ứng dây chuyền, nước mắt Pekora không ngừng chảy. Nó chảy, chảy mãi, tạo thành một dòng lũ cuồn cuộn, cuốn theo mọi thứ trên đường đi.

Chẳng mấy chốc, nước ngập đến bốn, năm mét, tới mức ngập cả thành phố.

Và như một phép màu kỳ diệu (hay một kịch bản mà chẳng ai ngờ tới), miếng cá kusuya mà Flare đang cầm trên tay - thứ vũ khí hủy diệt mùi hương ấy - cũng theo dòng nước mắt mà tự bơi về quê nhà của mình. Nó lướt nhẹ trên mặt nước, thỉnh thoảng vỗ vây, như thể đang tạm biệt thế giới loài người để trở về với biển cả.

\begin{figure}[H]
  \centering
  \begin{subfigure}[t]{0.45\textwidth}
    \centering
    \includegraphics[width=8cm]{../../../../Image/Book5/5-2l1.png}
    \caption{``Ở giữa thủ đô xa hoa tráng lệ\dots''}
  \end{subfigure}
  \quad
  \begin{subfigure}[t]{0.45\textwidth}
    \centering
    \includegraphics[width=8cm]{../../../../Image/Book5/5-2l2.png}
    \caption{``\dots\ con cá tung tăng bơi dưới lòng đường!''}
  \end{subfigure}
\end{figure}

Flare đứng nhìn dòng lũ, tay vẫn còn dang ra, mặt đầy vẻ bất lực: ``Tôi chỉ định cho cậu mượn củ cà rốt thôi mà\dots\ Sao thành lụt cả thủ đô thế này? Đây là phim thảm họa hay gì?''

Pekora vẫn khóc, nhưng lần này là khóc trong hạnh phúc – vừa sụt sùi vừa mỉm cười nhìn củ cà rốt trên đầu: ``Cảm ơn cậu, Flare\dots\ peko\dots''

Từ xa, cậu Tổng và Marine thò đầu ra khỏi cửa bếp, mặt tái xanh: ``Mấy người có thể giải thích tại sao văn phòng lại thành sông Hằng được không!? Bộ đang quay video ca nhạc dưới nước à!?''

Nhưng chẳng ai buồn trả lời. Vì ở khoảnh khắc ấy, một miếng cá khô đã trở về với đại dương, và một nàng Thỏ đã tìm lại được củ cà rốt của mình.

Kết thúc một ngày đầy mùi, đầy lệ, và đầy\dots\ những điều khó giải thích.

%==========================================TRUYỆN PHỤ=========================================
\omk

Sau một hồi vật lộn với thiên tai (nước mắt Pekora) và địa họa (cá thối bất tử), cuối cùng Flare và Pekora cũng quay trở về văn phòng. Cô nàng Yêu tinh vàng vừa bước vào cùng với nàng Thỏ tóc xanh, mắt cô vẫn còn hơi đỏ hoe sau trận khóc đến mức lụt cả thành phố. Trên tóc Pekora, củ cà rốt mới vẫn nằm yên vị, lấp lánh dưới ánh đèn.

``Bọn tôi về rồi đây\~{}\dots'' Flare cất giọng, định bụng tận hưởng bầu không khí trong lành, dễ chịu sau khi đã tống khứ được ``cỗ máy tận thế'' mang tên kusuya.

Nhưng chưa kịp thở phào, cả hai lập tức khựng lại.

Một mùi hôi khủng khiếp khác, Một mùi mới, lạ, nhưng không kém phần mãnh liệt; thứ này tấn công khứu giác của họ như một cơn sóng thần thứ hai. Pekora gắt lên, hai tay bịt mũi: ``Sao lại thối thế hả trời, peko!? Chẳng phải bọn mình vừa vứt con cá khô đi rồi sao? Hay nó đã sinh sản vô tính?''

Flare gật gù, mặt nhăn nhó: ``Tôi cảm tưởng thứ mùi này còn nặng hơn mùi kusuya ban nãy\dots\ Hôi quá! Giống như\dots\ có ai đó ngâm cá chết trong giấm lâu năm rồi mở nắp ra vậy.''

Trong phòng, Marine và Rushia đang cố gắng quạt cái mùi đó ra cửa sổ bằng hai chiếc quạt giấy nhỏ xíu, một nỗ lực vừa anh dũng vừa vô vọng. Mồ hôi lấm tấm trên trán họ, mắt thì nheo lại, miệng lẩm bẩm những câu gì đó không rõ. Noel thì vẫn đang nằm bất tỉnh ở góc phòng, trên mũi được nhét vài mảnh giấy ăn trông như một xác ướp Ai Cập phiên bản rẻ tiền.

Riêng Kuon vẫn ngồi bình thản trên ghế sô-pha, chân vắt chéo, tay cầm điện thoại đọc báo, mặt không hề tỏ vẻ khó chịu. Khói độc, mùi thối, xung quanh hỗn loạn. Cậu vẫn thư thái như một vị thiền sư giữa chiến trận.

Rushia vừa dùng chiếc quạt tay phẩy ra cửa sổ, vừa bịt mũi, giọng nghẹn lại: ``Tôi tưởng cái mùi cá thối này bay hết rồi chứ-nanodesu!? Sao còn nồng thêm thế này? Hay là con cá đó đã để lại linh hồn?''

Marine đáp lại, hợp lực với Rushia để đuổi cái mùi khó ngửi, nhưng mỗi phẩy quạt, mùi càng lan rộng: ``Thì tôi cũng có biết gì đâu chứ! Ặc\dots\ Mùi này\dots\ nồng quá\dots\ Giống như ai đó đã ủ một mẻ mắm mực (shiokara) trong tủ từ mấy tháng trước vậy!''

Nàng Pháp sư thảng thốt: ``Mắm mực á? Thứ đó còn khủng khiếp hơn cá khối gấp trăm lần!''

Thấy các đồng nghiệp phản ứng dữ dội, cậu quản lý chỉ nhún vai, thản nhiên nói, mắt vẫn không rời màn hình: ``Anh thấy mấy đứa cứ hơi quá. Bình thường như cân đường hộp sữa thôi. Hít một hơi sâu, cảm nhận, rồi sẽ quen.''

Lập tức, cả bốn cô gái của thế hệ thứ Ba đồng loạt quay sang nhìn cậu như thể cậu là một sinh vật ngoài hành tinh vừa đáp xuống trái đất. Mắt họ mở to, miệng há hốc, lòng đầy phẫn nộ.

Flare lên tiếng, giọng đầy bất lực: ``Sao anh có thể quen với mùi này được nhỉ? Tôi đang nghi ngờ khứu giác của anh có vấn đề.''

Pekora gật đầu lia lịa, mắt vẫn đỏ hoe: ``Có khi anh ấy là thành con peko nó rồi, peko! Chịu được mùi thối đến mức này chẳng khác gì siêu nhân.''

Rushia cũng không tha, giọng đầy buộc tội: ``Anh bị ấm đầu rồi sao!? Hay anh đã từng sống trong bãi rác nên miễn nhiễm rồi?''

Marine thì chỉ biết hét lên trong tuyệt vọng, tay vẫn phẩy quạt không ngừng: ``Có ai đó đuổi hết mùi thối này ra khỏi đây đi chứ!? Em không chịu nổi nữa! Mũi em sắp rụng rồi!''

Bất chợt, nàng Thỏ nhìn thấy một thứ gì đó lấp ló sau cánh cửa tủ gỗ ở góc phòng. Một cái bóng mờ mờ, hình trụ, màu tím đen, nằm im thít như một quả bom hẹn giờ đang chờ phát nổ.

``Mọi người ơi, nhìn nè!'' cô chỉ tay, giọng run run.

Cả bốn người dồn mắt theo hướng nàng Thỏ chỉ. Đó là một lọ dung dịch màu tím được đựng trong một chiếc lọ thủy tinh có nắp đậy kín, nhưng hình như nắp đã bị nới lỏng từ lúc nào. Một thứ ánh sáng huyền bí (hay đúng hơn là thứ hắc khí vô hình) tỏa ra từ khe nắp.

Họ từ từ tiến lại gần. Càng tới, cái bóng càng rõ hơn, nhưng mùi hôi cũng càng lúc càng nặng theo từng bước chân. Nó như một con quái vật vô hình đang mở rộng lãnh thổ, nuốt chửng mọi khứu giác trong tầm ảnh hưởng.

Flare đi đầu tiên, tay bịt mũi, giọng than thở đầy sợ hãi: ``Là cái gì thế\dots? Mùi quái gì mà ghê vậy!? Tôi chưa từng ngửi thấy thứ gì như thế này trong cả đời, kể cả khi vào thăm nhà máy xử lý rác!''

Rushia theo sát phía sau, mặt tái mét: ``Không lẽ lại là một con cá thối khác!? Hay là\dots\ natto? Tôi có linh cảm không lành về chuyện này lắm\dots''

Marine ở ngay bên cạnh, tỏ vẻ gắt gỏng nhưng cũng không giấu được sự tò mò: ``Tôi chỉ mang có mỗi một miếng cá khô thôi mà! Cái này chắc chắn không phải hàng của tôi!''

Pekora là người cuối cùng, nuốt nước bọt, giọng nghẹn lại: ``Không lẽ trong tủ có\dots\ ác quỷ hôi thối à, peko!? Hay là xác ướp của một con cá mập khô lâu năm?''

Tất cả bọn họ nhìn nhau, chẳng ai dám mở tủ ra. Họ đứng thành vòng tròn xung quanh chiếc tủ, mắt dán vào cánh cửa, như thể đang chuẩn bị đối mặt với một con quái vật sắp nhảy ra. Flare cuối cùng miễn cưỡng tiến lên trước, một tay che mũi, tay kia run run đặt lên nắm cửa.

Cô hít một hơi thật sâu (và lập tức hối hận vì đã hít), rồi quay về phía mọi người, giọng như sắp khóc: ``Được rồi, tôi sẽ mở\dots\ nhưng nếu có gì nhảy ra thì đừng trách tôi đập nó nha! Tôi đang cầm sẵn dép lào đây!''

Pekora thấy vậy, thốt lên đầy hoảng sợ: ``Chờ đã! Hay là ta mở hé một chút rồi nhìn vào trước—có khi chỉ là bình nước lọc bị hỏng thôi mà, peko!''

Nhưng đã quá muộn.

Không cần mở nhẹ nhàng, không cần chuẩn bị tâm lý. Cánh cửa tủ tự bật tung ra như thể thứ gì đó bên trong đã không thể chờ đợi lâu hơn nữa, một lực vô hình bung ra, như một quả bom nổ chậm sau mấy tháng ủ mình.

Một luồng khí nồng nặc, đặc quánh, mạnh đến mức gần như có thể nhìn thấy bằng mắt thường. Một thứ khí màu tím nhạt, bốc lên như sương mù lập tức tràn ra, xộc thẳng vào mặt cả bốn người.

Rushia ho sặc sụa, bị ngợp với mùi hương khó tả, mắt trào nước mắt: ``\textbf{C-Cái quái gì thế này-nanodesu!? Phổi em như bị nấm mốc xâm chiếm rồi!}''

Marine lùi lại, cố gắng tránh cái mùi hương đó trong tuyệt vọng, nhưng nó bủa vây khắp nơi, không có lối thoát: ``\textbf{Trời ơi, mùi này là tận thế luôn rồi! Ai đó cứu tôi với!! Em xin lỗi, em sẽ không bao giờ mang cá khô vào công ty nữa!}''

Flare cũng không khá hơn. Cô choáng váng, lảo đảo khắp phòng như một con tàu trong cơn bão, tay bám vào tường để khỏi ngã: ``Khốn kiếp\dots\ Không ngờ lại nặng đến mức này\dots\ Tôi tưởng mình đã từng ngửi qua mọi thứ rồi chứ\dots''

Pekora thì trợn ngược mắt, hai tay bóp cổ, nước mắt giàn giụa, nghĩ rằng mình sắp chầu trời đến nơi: ``T-Tôi không ổn\dots\ peko\dots\ Củ cà rốt trên đầu tôi\dots\ đang rụng hết\dots\ peko\dots''

Và rồi, từng người một, như những quân cờ domino đổ xuống, cả bốn cô gái đồng loạt lăn đùng ra sàn. Họ nằm sõng soài bên cạnh nàng Hiệp sĩ Noel vẫn đang bất tỉnh, mắt quay cuồng, tay chân co quắp, không khác gì những chiến binh đã bại trận trước kẻ địch vô hình nhưng cực kỳ đáng sợ.

Flare nằm nghiêng, thều thào: ``Đây\dots\ là kết cục của những kẻ\dots\ tò mò sao\dots?''

Marine úp mặt xuống sàn, giọng tắt dần: ``Em\dots\ hối hận\dots\ quá\dots''

Rushia gục đầu lên chân Marine, miệng lẩm bẩm: ``Nanodesu\dots\ nanodesu\dots\ em về với Chúa đây\dots''

Pekora nằm ở giữa, hai tay vẫn ôm chặt củ cà rốt trên đầu, có lẽ là vật duy nhất cô định đưa lên thiên đường cùng.

``Có chuyện gì mà mấy đứa cứ quay cuồng như mấy con rối thế?''

Cậu Tổng cảm thấy lạ, bước lại nhìn bốn cô nàng đã bị hạ đo ván. Xung quanh họ, bầu không khí vẫn còn lưu lại một lớp sương mù mờ ảo, màu tím nhạt, phảng phất mùi kinh dị.

Cậu liếc vào bên trong chiếc tủ, rồi bất chợt, ký ức bỗng tràn về. Một ký ức xa xôi, từ cái thời mà cậu còn\dots\ thích ăn bún đậu.

``À\dots\ cái lọ mắm tôm anh để đây từ mấy tháng trước này! Hôm đó mang đi định ăn bún đậu cùng mấy đứa bạn, nhưng gặp mưa to quá nên không đi được, về quên luôn trong tủ. Hóa ra là nó ở đây!''

Cậu nhìn xuống những nạn nhân đang nằm bất tỉnh nhân sự, đặc biệt là Flare, ``lính mới'' vừa gia nhập đã phải chịu trận ngay ngày đầu. Trông cô chẳng khác gì đang trải qua một cơn ác mộng kinh hoàng, môi mấp máy những gì không rõ.

Chứng kiến cảnh tượng khó đỡ này, cậu Tổng chỉ biết lắc đầu, khẽ cười: ``Xin lỗi mấy đứa nha. Anh quên mất. Lần sau sẽ để ý.''

Cái lọ mắm tôm để quên hàng tháng trời đó vẫn vô tư nằm yên trong tủ, nắp đã mở hé, tiếp tục tỏa ra thứ sức mạnh vô song, thứ sức mạnh đã hạ gục toàn bộ thế hệ thứ Ba chỉ trong vòng chưa đầy ba phút.

Ngoài cửa sổ, gió vẫn thổi, trời vẫn sáng. Và từ xa, có thể nghe thấy tiếng cậu Tổng thì thầm: ``Mà\dots\ để anh đóng nắp lại, cất kỹ. Biết đâu lúc thiếu gia vị lại có mà dùng.''

\end{document}

